% SHARD-ID: AQM-65-QSA-VOCABULARY
% SHARD-TITLE: Quantum-System Association Vocabulary
% SHARD-SUMMARY: Fixes the programme-local vocabulary for quantum-system associations before comparison shards begin.
% SHARD-SUMMARY: Separates association workflows from any canonical functorial quantisation claim.
% SHARD-SUMMARY: Records the evidence rule for later agreement, equivalence, inequivalence, and degeneracy claims.
% SHARD-KEYWORDS: quantum-system association, vocabulary, Quant1, Quant2, equivalence, degeneracy, kinematics

\section{Quantum-System Association Vocabulary}
\label{sec:quantum-system-association-vocabulary}

\subsection{Status and Local Anchors}

This is a planning and convention shard only.  It introduces no substantive
equivalence theorem, no sourced mathematical classification, and no new
representation-theoretic assertion.  Its local anchors are
\texttt{CONVENTIONS.md} convention \texttt{(ap)}, the meta-plan shard
\texttt{AQM-64-QUANTUM-SYSTEM-ASSOCIATION-META-PLAN}
(\S\ref{sec:quantum-system-association-meta-plan}), and Step 01 of
\texttt{docs/quantum\_system\_associations/PLAN.md}.

\subsection{Association, Not Canonical Quantisation}

In this programme, a \emph{quantum-system association} is a recorded workflow
which starts from specified structure on an object \(X\), adds explicitly
named choices, and returns kinematical quantum data.  Possible output slots
include a Hilbert carrier, observable algebra, Weyl--Heisenberg projective
representation, Fock sector, or stabilizer subspace, but the chosen slot must
be named in the relevant shard.

This language deliberately does not assert a canonical functorial
quantisation.  A later shard may study functorial behaviour for a specified
workflow only after recording the objects, morphisms, choices, and evidence.
Until then, the word ``association'' means a reproducible construction
proposal with an evidence status, not a theorem that one universal
quantisation functor exists.

\subsection{Local Shorthand}

The following terms are local programme labels.

\begin{description}
\item[\texttt{Quant1}, \texttt{Quant2}.]
Placeholders for the first and second association workflows in a comparison.
They carry no fixed mathematical content until a shard expands them into
explicit inputs, choices, and outputs.

\item[field association.]
An association whose planned kinematical output is organized as field-like
data over points, opens, supports, sections, or labels.  The term is an
umbrella for comparison planning, not a claim that all such workflows are
instances of one established construction.

\item[single-particle sector.]
A comparison label for the part of a workflow in which the carrier represents
one alternative among sites or modes rather than independent coexistence of
all sites.  A later shard must say whether this means a direct sum, a selected
Fock sector, or another explicitly recorded operation.

\item[association equivalence.]
An evidence status for two association workflows after a concrete object
\(X\), all needed choices, and a checkable comparison map or invariant have
been recorded.  Without that data, ``equivalence'' remains a question or
proposal.

\item[degenerate case.]
A case in which a workflow collapses, forgets the intended structure, becomes
too small or too large for the intended comparison, or reduces to an earlier
case.  Calling a case degenerate requires the same local evidence discipline
as any other comparison status.
\end{description}

\subsection{Dynamics Deferred}

The kinematical output of a quantum-system association does not include
dynamics by default.  Hamiltonians, actions, time evolutions, Witten-index
statements, and dynamical selection principles are later choices unless a
future convention and shard explicitly add them.

\subsection{Evidence Rule for Later Comparisons}

Every future comparison word such as ``agrees'', ``reduces to'', ``is
equivalent to'', ``is inequivalent to'', or ``degenerates'' must cite a local
source, local derivation, test, run artifact, or explicit gap label.  The
comparison must name the object \(X\), the structure placed on \(X\), the
choices made by each workflow, the kinematical outputs being compared, and the
map or invariant used for the comparison.  Otherwise the statement must remain
an open question in the report.
